\documentclass[9pt,landscape,a4paper]{article}
\usepackage[utf8]{inputenc}
\usepackage[T1]{fontenc}
\usepackage{geometry}
\usepackage{multicol}
\usepackage{amsmath,amssymb}
\usepackage{enumitem}
\usepackage{graphicx}
\usepackage{xcolor}
\usepackage{listings}
\usepackage{booktabs}
\usepackage{titlesec}
\usepackage{fancyhdr}
\usepackage{hyperref}
\usepackage{CJKutf8}

% Page layout
\geometry{top=0.7cm, bottom=0.7cm, left=0.7cm, right=0.7cm}
\pagestyle{empty}
\setlength{\parindent}{0pt}
\setlength{\parskip}{1.5pt}
\setlength{\columnsep}{7pt}

% Section formatting
\titleformat{\section}{\normalsize\bfseries\color{blue!70!black}}{}{0em}{}[\titlerule]
\titleformat{\subsection}{\small\bfseries\color{blue!50!black}}{}{0em}{}
\titlespacing{\section}{0pt}{4pt}{2pt}
\titlespacing{\subsection}{0pt}{3pt}{1pt}

% Code listing style
\lstset{
    language=Python,
    basicstyle=\ttfamily\fontsize{6.5}{7.5}\selectfont,
    keywordstyle=\color{blue!70!black}\bfseries,
    stringstyle=\color{red!60!black},
    commentstyle=\color{green!50!black}\itshape,
    numberstyle=\tiny\color{gray},
    backgroundcolor=\color{gray!8},
    frame=single,
    framerule=0.3pt,
    rulecolor=\color{gray!50},
    breaklines=true,
    breakatwhitespace=false,
    tabsize=2,
    showstringspaces=false,
    aboveskip=2pt,
    belowskip=2pt,
    xleftmargin=2pt,
    xrightmargin=2pt,
}

\newcommand{\keybox}[1]{%
    \begin{center}
    \fcolorbox{blue!50!black}{blue!5}{%
        \parbox{0.95\linewidth}{\small #1}%
    }
    \end{center}
}

\newcommand{\zhint}[1]{{\scriptsize\color{gray!70!black}#1}}

\begin{document}
\begin{CJK}{UTF8}{gbsn}

\begin{center}
{\large\bfseries\color{blue!70!black} Chapter 12 -- Case Studies (IV): Game Problems}\\
{\footnotesize IERG4320/IEMS5726 Data Science in Practice | Cheat Sheet}
\end{center}

\begin{multicols}{3}

% ==================== SECTION 1 ====================
\section{Game AI Overview}
\zhint{游戏AI：用搜索和学习找到最优策略}

\textbf{Games \& AI}: well-defined rules (env), actions, goals (rewards). AI serves as opponents, teammates, commentators, story guides.

\textbf{Search techniques}: A* search, alpha-beta pruning (minimax), fuzzy logic, genetic algorithm.\\
\zhint{搜索技术：A*、Alpha-Beta剪枝、模糊逻辑、遗传算法}

\textbf{Rule-based}: if winning move $\to$ win; if opponent winning $\to$ block; else promote or random. Simple but limited for complex situations.

% ==================== SECTION 2 ====================
\section{Blackjack Environment}
\zhint{21点：不超过21点的前提下赢过庄家}

\textbf{Objective}: beat dealer without going over 21.\\
Cards: 2--10 face value; J/Q/K = 10; Ace = 1 or 11.

\textbf{State}: tuple (player\_sum, dealer\_showing, usable\_ace).\\
E.g. (15, 10, False): player has 15, dealer shows 10, no usable Ace.\\
\textbf{Actions}: Hit (take card) or Stand (stop).\\
\textbf{Reward}: +1 win, --1 loss, 0 draw (sparse, end of episode only).\\
\textbf{Dealer policy}: must hit until sum $\geq$ 17.\\
\zhint{庄家必须在<17时继续要牌}

\textbf{Basic strategy}: win 42.2\%, loss 49.1\%, tie 8.5\%.

% ==================== SECTION 3 ====================
\section{D3QN (Dueling Double DQN)}
\zhint{D3QN：双重DQN减少过估计 + 对决架构分离状态价值和优势}

\subsection{Double DQN}
\zhint{双重DQN：用在线网络选动作，目标网络评估}\\
Fixes overestimation bias in standard DQN.\\
Standard: $y = r + \gamma \max_{a'} Q_{tar}(s', a'; \theta^-)$\\
Double: $y = r + \gamma \, Q_{tar}(s', \arg\max_{a'} Q(s', a'; \theta); \theta^-)$\\
\zhint{用$\theta$选动作，用$\theta^-$评估价值}

\subsection{Dueling Architecture}
\zhint{对决架构：分成价值流V(s)和优势流A(s,a)}\\
Splits Q-network into two streams:\\
\textbf{Value} $V(s)$: how good is this state?\\
\textbf{Advantage} $A(s,a)$: how much better is action $a$ than average?\\
$Q(s,a) = V(s) + \big(A(s,a) - \text{mean}_a A(s,a)\big)$\\
E.g. state (20,2,False): $V(s)$ high (great position), $A(\text{STAND})$ positive, $A(\text{HIT})$ very negative.

\textbf{Loss}: $\mathbb{E}[(y - Q(s,a;\theta))^2]$ where $\theta^-$ updated periodically.\\
\textbf{Result}: win rate $\approx$ 43--44\% after 10k games.

% ==================== SECTION 4 ====================
\section{A3C (Async Advantage Actor-Critic)}
\zhint{A3C：多个并行worker异步更新全局网络}

\textbf{vs DQN}: supports continuous actions, trains faster, avoids maximization bias.

\subsection{Key Components}
\begin{itemize}[leftmargin=8pt,itemsep=0pt,topsep=0pt]
\item \textbf{Asynchronous}: multiple workers explore in parallel, decorrelate data naturally \zhint{并行探索，自然去相关}
\item \textbf{Actor}: learns policy $\pi(a|s)$ \zhint{策略网络}
\item \textbf{Critic}: learns value $V(s)$ \zhint{价值网络}
\item \textbf{Advantage}: $A(s,a) = Q(s,a) - V(s)$\\
Positive $\to$ reinforce; Negative $\to$ discourage \zhint{正=强化，负=抑制}
\end{itemize}

\subsection{Training Pipeline}
\begin{enumerate}[leftmargin=10pt,itemsep=0pt,topsep=0pt]
\item Workers collect (s, a, r) sequences independently
\item Compute n-step returns $\to$ estimate advantage
\item Calculate actor + critic gradients
\item Asynchronously update global network
\end{enumerate}
\zhint{n-step回报：短n=高偏差低方差; 长n=低偏差高方差}\\
\textbf{Result}: win rate $\approx$ 38--40\% (tends to always STAND).

% ==================== SECTION 5 ====================
\section{AlphaZero / AlphaBlackjack}
\zhint{AlphaZero：DeepMind开发，只用规则+自我对弈学习}

Self-play only: given rules, plays millions of games against itself.

\subsection{MCTS (Monte Carlo Tree Search)}
\zhint{蒙特卡洛树搜索：选择→扩展→模拟→反向传播}
\begin{enumerate}[leftmargin=10pt,itemsep=0pt,topsep=0pt]
\item \textbf{Selection}: pick node via $Q + U$ (exploit + explore)
\item \textbf{Expansion}: add new move to tree
\item \textbf{Simulation}: play randomly to end (rollout)
\item \textbf{Backpropagation}: update stats along path
\end{enumerate}

\subsection{Dual-Headed Network}
\textbf{Policy head} (Actor): outputs $[P(\text{STAND}), P(\text{HIT})]$.\\
\textbf{Value head} (Critic): scalar $\in [-1, +1]$.\\
Training: policy gradient (actor) + MSE (critic).

\subsection{Self-Play Training}
Early: learn basics (don't hit on 20).\\
Middle: nuanced strategy (16 vs dealer's 7).\\
Late: master complex situations (soft hands).\\
Each iteration plays against previous version.\\
\textbf{Result}: win rate $\approx$ 38--40\%.

% ==================== SECTION 6 ====================
\section{Single Agent for Multiple Games}
\zhint{单智能体玩多个游戏：核心挑战是灾难性遗忘}

\textbf{Problem}: catastrophic forgetting -- learning Game B overwrites weights crucial for Game A.\\
\zhint{灾难性遗忘：学新游戏会覆盖旧游戏的重要权重}

\subsection{EWC (Elastic Weight Consolidation)}
\zhint{弹性权重固化：惩罚对旧任务重要参数的大改动}\\
Penalizes large changes to parameters important for previous tasks; allows flexibility for unimportant ones.\\
\textbf{Process}: train Game A $\to$ train Game B with EWC (important weights ``elasticated'' toward A's solution) $\to$ repeat for Game C (EWC terms for both A and B).\\
\textbf{Target games}: Blackjack (simple), Leduc Hold'em (imperfect info, bluffing), Hearts (long-term strategy).

% ==================== SECTION 7 ====================
\section{3D Model Generation}
\zhint{3D模型生成：从2D图像或文本生成3D模型}

\textbf{3D representations}: mesh (vertices + faces), depth map, voxel grid, point cloud, implicit surface.\\
\textbf{Triangle mesh}: standard for graphics, adaptive, renderable, less memory. 8 vertices + 12 faces = cube.
\begin{lstlisting}
import trimesh, numpy as np
vertices = np.array([
  [0,0,0],[1,0,0],[1,1,0],[0,1,0], # bottom
  [0,0,1],[1,0,1],[1,1,1],[0,1,1]  # top
])
faces = np.array([
  [0,1,2],[0,2,3], [4,5,6],[4,6,7], # bot/top
  [0,4,7],[0,7,3], [1,5,6],[1,6,2], # L/R
  [0,1,5],[0,5,4], [3,2,6],[3,6,7]  # F/B
])
mesh = trimesh.Trimesh(vertices, faces)
mesh.show()
\end{lstlisting}

\subsection{TripoSR (2D $\to$ 3D)}
\zhint{TripoSR：单张图片重建3D模型}\\
Input: single RGB image. Output: 3D mesh.\\
Pipeline: DINOv1 encoder $\to$ Transformer decoder $\to$ triplane NeRF $\to$ marching cubes $\to$ mesh.\\
No camera params needed (learned during training).\\
\textbf{Improvements}: 40-channel triplane, BCE mask loss (reduce floaters), local 128$\times$128 patch rendering.\\
\textbf{Metrics}: Chamfer Distance (lower=better), F-Score (higher=better).

\subsection{Shap-E (Text $\to$ 3D)}
\zhint{Shap-E：文本/图像生成3D隐式神经表示，约13秒}\\
Generates INR (Implicit Neural Representation) parameters via diffusion.\\
\textbf{Two-stage}: (1) text $\to$ 2D image (image diffusion); (2) image $\to$ 3D point cloud (point cloud diffusion conditioned on CLIP embedding).\\
\textbf{Upsampler}: low-res $\to$ high-res point cloud $\to$ Poisson surface reconstruction $\to$ mesh.\\
Speed: $\sim$13s per asset on V100 GPU.\\
\textbf{Limitations}: lower fidelity, floater artifacts, topological errors, bottlenecked by 2D model quality.

% ==================== SECTION 8 ====================
\section{Robot Learning}
\zhint{机器人学习：构建感知-动作循环}

Robot vision: embodied, active, situated.\\
Formulated as RL: perception $\to$ action $\to$ feedback loop.

\subsection{Transfer Learning}
\zhint{迁移学习：用预训练模型适应新任务}
\begin{itemize}[leftmargin=8pt,itemsep=0pt,topsep=0pt]
\item \textbf{Fixed feature extractor}: freeze pre-trained layers, train new classifier
\item \textbf{Fine-tuning}: continue backprop on pre-trained weights
\item \textbf{When}: small+similar $\to$ fix features; large+similar $\to$ fine-tune all; small+different $\to$ earlier layers; large+different $\to$ fine-tune all
\end{itemize}

\subsection{Robot Learning Methods}
\zhint{机器人学习方法}
\begin{itemize}[leftmargin=8pt,itemsep=0pt,topsep=0pt]
\item \textbf{Imitation Learning}: mimic expert demonstrations \zhint{模仿学习}
\item \textbf{Behavior Cloning}: supervised (state $\to$ action) \zhint{行为克隆}
\item \textbf{Inverse RL}: infer expert's reward function, then optimize \zhint{逆强化学习}
\item \textbf{Diffusion Policies}: generate multi-modal action sequences via diffusion models \zhint{扩散策略}
\end{itemize}

\keybox{
\textbf{Ch12 Summary -- 3 Perspectives}:\\
\textbf{Learn to play}: D3QN, A3C, AlphaZero, EWC\\
\textbf{Learn to design}: TripoSR (2D$\to$3D), Shap-E (text$\to$3D)\\
\textbf{Learn to walk}: Transfer learning, imitation, behavior cloning, inverse RL, diffusion policies
}

\end{multicols}
\end{CJK}
\end{document}
